\chapter{INTRODUCTION}

\section{Background of the Study}

Understanding how users navigate through digital systems---such as e-commerce websites, mobile applications, and customer support platforms---is essential for improving conversion rates, user satisfaction, customer retention, and operational efficiency. User interaction data captures detailed behavioral patterns that reveal how users navigate through different system components. Customer journeys, represented as sequences of interaction states (e.g., \emph{Landing $\rightarrow$ Browse $\rightarrow$ Product $\rightarrow$ Cart $\rightarrow$ Checkout}), provide a structured framework for analyzing these behaviors.

Modern digital systems generate millions of user sessions across thousands of interaction states. When aggregated, these journeys form large directed graphs whose scale and complexity necessitate efficient computational analysis grounded in graph theory and data mining.

\section{Statement of the Problem}

As the volume and complexity of user interaction data continue to grow, identifying meaningful and optimal customer journeys becomes increasingly challenging. From a business perspective, customer journey dynamics directly influence conversion rates, revenue generation, customer retention, and user satisfaction. Even small improvements in the probability of reaching a conversion state can result in significant financial gains at scale.

Na\"{i}ve analytical methods, such as exhaustive path enumeration or simple frequency-based techniques, do not scale. In large graphs, the number of possible simple paths increases combinatorially, rendering exhaustive enumeration computationally infeasible. Furthermore, directly multiplying many small transition probabilities along a path leads to numerical underflow---a form of arithmetic precision loss that makes na\"{i}ve probability computation unreliable for longer paths.

The key challenge is therefore to identify optimal customer journeys efficiently and with numerical stability at scale. Core tasks include finding highest-probability conversion paths and analyzing scalability across large graphs. Each task poses distinct challenges in graph traversal, optimization, and algorithm design that require careful data structure selection.

Maximizing path probability is equivalent to optimizing for the most likely conversion sequence. In an e-commerce context, the highest-probability path from a landing page to checkout represents the conversion funnel that customers naturally follow. Identifying this path allows businesses to reduce friction along the dominant route, prioritize UI investments on the states that carry the most traffic, and design nudges that align with observed user behavior rather than assumed behavior. The term ``optimization'' in the project title therefore refers not only to algorithmic speedup but also to the actionable business insight that optimal-path discovery enables.

\section{Objectives}

The objective of this study is to design and implement an efficient algorithmic framework for customer journey path optimization grounded in graph theory and classical algorithm design principles. Specifically, the project models customer journey data as a directed weighted graph, where nodes represent interaction states and edges represent transition probabilities between states.

The primary focus is to compute optimal conversion paths under probabilistic constraints. To achieve this, transition probabilities are transformed into additive non-negative edge weights using the negative logarithm, enabling the application of Dijkstra's shortest-path algorithm for maximum-probability path computation. This transformation also addresses the numerical underflow problem by converting multiplicative probability products into additive sums. The project implements Dijkstra's algorithm as a baseline approach and proposes an optimized probability-pruned search strategy that reduces the exploration of low-probability paths to improve computational efficiency.

In addition to correctness, the project formally analyzes the theoretical time and space complexity of both the baseline and optimized algorithms. Empirical performance evaluations have been conducted on graphs of varying sizes and densities to examine scalability, runtime behavior, and trade-offs between computational efficiency and path optimality. Through these analyses, the project demonstrates how careful algorithm design and data structure selection influence performance in large-scale navigation graphs.


\noindent\textbf{Contribution:} This study presents a \textbf{well-motivated, controlled pruning framework} for probabilistic shortest-path computation, grounded in classical algorithm design and supported by comprehensive empirical evaluation. The pruning variant is deterministic, requires no domain-specific heuristic (unlike A* search), and provably converges to the globally optimal solution as $\tau \to 0$ (unlike rank-based methods such as beam search, which provide no optimality guarantee). The main contributions are: (1)~formal characterization of an ``all-or-nothing'' correctness property (Theorem~\ref{thm:all_or_nothing}), stating that the algorithm returns either the exact optimal path or reports no path found with no sub-optimal intermediate results; (2)~identification and empirical characterization of the \emph{critical pruning threshold} $\tau_c$, a narrow operational range where speedup is maximized while maintaining optimality; and (3)~systematic evaluation framework spanning synthetic and real-world datasets, demonstrating that the approach scales to graphs with tens of thousands of nodes while preserving solution quality.

\medskip
\noindent\textbf{Novelty and Prior Work:} To the best of our knowledge, this is the first report of a deterministic, parameter-controlled probability-pruned Dijkstra algorithm with a formal all-or-nothing property and an adaptive-$\tau$ regime for real-world graphs. While pruning, heuristics, and thresholding are well-studied in shortest-path literature, no prior work combines these elements with formal guarantees and empirical validation in the context of probabilistic customer journey graphs. Our findings and methodology are not reported elsewhere in the literature.

\section{Research Hypothesis}

The following hypothesis guides the experimental evaluation:

\begin{quote}
Probability-Pruned Dijkstra provides meaningful speedup over baseline Dijkstra while preserving path optimality at conservative $\tau$ values.
\end{quote}

\noindent The experimental evidence supporting this hypothesis is presented in Chapter~\ref{chap:results}; a summary of findings is provided in Chapter~\ref{chap:conclusion}.

\section{Scope and Applicability}

The framework developed in this study is applicable to domains where user interactions can be modeled as probabilistic directed graphs, including e-commerce conversion funnels, mobile application navigation flows, and customer support routing systems. Beyond industrial applications, the algorithmic contributions---particularly the formal correctness guarantees and the characterization of the reachability--speed trade-off---are relevant to researchers investigating threshold-based pruning strategies in shortest-path computation.
